\chapter{Introduction}
\label{sec:introduction}

% Source notes: Project/docs/implementation_plan.md §1–2 (research questions), checkpoint_2026-10-10.md (results),
% thesis_notes.md (locality pillar; Phase 0 prefill/decode timings), appendix.tex (example answer rt-7784).

\section{Motivation}
\label{sec:intro-motivation}

\subsection{Retrieval-augmented generation on the user's device}
\label{sec:intro-local}

Retrieval-augmented generation has become the standard way to let a language model answer questions about documents it
was never trained on. A retriever selects the passages that are relevant to the question, and the model writes its answer
with those passages in its context~\cite{lewis2020rag}. The same recipe serves a company searching its internal reports,
a lawyer reading a contract and a student questioning their own notes, and in each case the value of the answer depends
on its staying faithful to the documents.

Until recently such a pipeline required a hosted model, reached over the network. This is no longer the case. Compact
open-weight models and efficient runtimes now run on consumer hardware~\cite{llamacpp,mlx2023}, and platform vendors ship
models of a few billion parameters directly on their devices~\cite{gunter2024apple}. The gap with the best closed models
is small and shrinking: models that fit on a single consumer graphics card reach the performance of the frontier of the
time within six to twelve months~\cite{epoch2025consumergpu}, and the lead of the best closed model over the best open one
on a popular leaderboard fell from 8\,\% to under 2\,\% in a single year~\cite{stanford2025aiindex}. The whole pipeline, from retrieval
to answer, can therefore stay on the user's own machine. This matters for four reasons. When the documents are private
or regulated, they never leave the device; privacy becomes a technical property of the system rather than a promise of a
provider, which the user cannot verify. A recent dispute over credit between OpenAI and two mathematicians, who had
written unpublished drafts inside one of the company's products, showed how real the question is: the provider itself
stated that it could not rule out that data derived from that usage had helped improve its
models~\cite{rits2026navierstokes}. When there is no network, the assistant keeps working. When the behaviour of
the system must stay stable, the weights change only if the user changes them, instead of following a provider's
updates. Finally, the cost becomes that of the hardware and its energy, instead of a price per token. Local models also
expose their logits, which hosted interfaces often hide or limit, and this property turns out to be essential for the
approach studied here.

Locality is a constraint as well as an advantage. A laptop has limited memory, limited bandwidth and a battery, and every
component of the assistant must fit within them at the same time. This thesis treats locality as a requirement rather
than as an afterthought: every step of the work, from building the dataset to training, evaluation and serving, runs on
one consumer laptop with 16\,GB of memory, and no step depends on an external service. Because locality is the backbone
of the work, Chapter~\ref{sec:local} is devoted entirely to it: why it matters, which models and hardware make it
possible, and what it costs.

\subsection{Hallucination and the need for verification}
\label{sec:intro-hallucination}

Retrieval does not remove hallucination. Even with the right passages in context, generated answers still contain
sentences that the passages do not support, or that contradict them~\cite{niu2024ragtruth,huang2025hallucination}. An
example from the data used in this thesis shows how subtle the problem can be. Asked to describe a restaurant from a
structured record, a model writes six fluent sentences. Four of them are correct. One states that the restaurant is open
seven days a week, whereas the record shows that it is closed on Mondays: the sentence contradicts the source. Another
states that the restaurant has a rating of 4.5 stars on Google, whereas the record gives the rating without naming any
source: the sentence adds a detail that the source does not mention (Appendix~\ref{app:examples}).

Two observations follow from this example. First, the error sits in a single sentence of an otherwise correct answer, so
a verdict on the whole answer is of little help: the user needs to know which sentence to distrust. Second, the two
erroneous sentences are wrong in different ways. A contradiction is a factual error with respect to the source, whereas
an unsupported addition may be harmless background or an invention, and an application may want to treat the two
differently. A useful verifier therefore decides, for every sentence, whether the document supports it, contradicts it or
does not mention it. Since no user can be expected to check every sentence against its source, the verifier must run on
every answer, automatically, and on the same machine that generated it.

\subsection{The cost of checking an answer on a laptop}
\label{sec:intro-cost}

A verifier that runs next to the generator competes with it for the same memory, time and battery. The most common
solution, asking a language model to judge the answer, is expensive precisely where a laptop is weakest. A judge writes
its verdicts token by token, and on local hardware each generated token costs far more than each token it reads. Reading
the prompt is limited by computation, which a laptop GPU handles well, whereas writing the output is limited by memory
bandwidth, because every new token must read all the weights~\cite{pope2023scaling}. On the laptop used in this thesis,
reading one token takes about 0.37 milliseconds for a model of 0.5B parameters, while writing one takes about 5.9, sixteen
times more, and the asymmetry holds on every class of local hardware, from phones to dedicated AI machines
(Chapter~\ref{sec:local}). A judge that explains its reasoning before each verdict multiplies this cost by the length of its
explanation, and the cost grows again with the number of sentences in the answer. Small judges also break the required
output format, and the confidence they express is rarely usable as a probability.

Encoder-based verifiers avoid generation altogether. Models such as AlignScore, MiniCheck and HHEM read a pair of document
and sentence and return a score of support~\cite{zha2023alignscore,tang2024minicheck,vectara2024hhem}. They are small and
fast, but they give one fixed binary score per pair: they do not distinguish a contradiction from an omission, they read
each sentence without the rest of the answer, and adapting them to a new decision requires a new model.

\section{System One models}
\label{sec:intro-systemone}

A different way to use a language model for decisions has recently attracted attention under the name of \emph{System
One models}~\cite{typesafe2026jev}. The name evokes the fast, intuitive mode of thinking, as opposed to deliberate
reasoning. Instead of generating an answer, such a model reads the question and the context once and returns a
probability for each admissible answer, for example by reading the logits of the option letters at the end of the
prompt~\cite{semif,rizzoflow}. Three properties follow directly. No text is generated, so the decision costs only the
reading of the input. The output is always one of the admissible answers, so it is always valid. Many questions about the
same context can share a single pass over it, so each additional question costs little.

The idea fits verification well. An answer of $K$ sentences poses $K$ closed questions about the same document, and each
of them has three admissible answers. A System One can answer all of them while reading the document only once, and the
probability it returns for each sentence can be calibrated and compared with a threshold chosen by the application.

The field is young and moves quickly. After the first hosted model, open implementations
followed~\cite{semif,rizzoflow}, and in 2026 the Hugging Face hub introduced a tag for these \emph{decision models} which
by October 2026 already grouped about 1{,}500 models~\cite{hf2026decisionmodels}. Established organisations released their own,
such as Liquid AI's d1-3B and Cloudflare's Clef~\cite{liquid2026d1,cloudflare2026clef}. The evidence, however, lags
behind the releases. Most of it comes from the authors of each system, independent evaluations are recent and often less
favourable~\cite{deusser2026jev,nandakishore2026laya}, and a common leaderboard finds that general-purpose decision models
with fewer than one billion parameters barely exceed the trivial baseline on hallucination detection in retrieval-augmented
answers~\cite{decisionindex2026}. All the open implementations, moreover, are general-purpose systems that answer any
typed question.

Several questions therefore remain open. It is not known whether the approach is worth it for one specific task, how a
small model specialised in that task compares with a general-purpose one, with generative judges and with dedicated
verifiers, and where any advantage actually comes from: from the specialisation, from the way the decision is read, or
simply from the size of the model.

\section{Approach and research questions}
\label{sec:intro-rq}

This thesis designs, trains and evaluates GroundOne, a compact System One specialised in one decision: for every sentence
of a RAG answer, whether the retrieved document supports it, contradicts it or does not mention it. The answer is written
with a short marker after each sentence, and the verdicts of all sentences are read in one forward pass, from the logits
of three answer letters at the markers, without generating any text. Qwen2.5 models with 0.5B and 3B parameters are
fine-tuned with LoRA on human-labelled answers from RAGTruth~\cite{qwen2024qwen25,hu2022lora,niu2024ragtruth}, and their
probabilities are calibrated on development data. Three test sets were built and frozen before any model was trained, so
that no design choice could be tuned on them. Five research questions guide the work.

\begin{description}
\item[RQ1, specialisation.] Can a fine-tuned System One of 0.5B or 3B parameters match larger zero-shot models and the
specialised encoders on sentence-level groundedness, that is, on deciding whether each sentence is faithful to the
document?
\item[RQ2, readout versus generation.] What does the System One mechanism, the reading of the decision from the logits,
contribute at the same model size? The comparison is twofold: against generative judges that write their verdicts, and
against the same model deciding by the token it would generate. The question also asks how the difference grows with the
number of decisions per answer.
\item[RQ3, calibration.] Are the probabilities usable as confidence, to set thresholds and to defer uncertain decisions,
and how do they compare with the probabilities that a generative judge can provide?
\item[RQ4, cost on the device.] Does the verifier fit within the memory, time and energy budget of a 16\,GB laptop, and
how does its cost compare with that of a judge running on the same machine?
\item[RQ5, beyond a binary score.] What does a three-way decision for every sentence offer that a binary verifier does
not?
\end{description}

\section{Scope}
\label{sec:intro-scope}

The thesis studies the verification step in isolation. The retrieved document is given, so the quality of retrieval is
not evaluated, and an answer is judged only against the document it was generated from, not against world knowledge.
GroundOne flags sentences but does not correct them, and it returns verdicts and probabilities rather than written
explanations. The data are in English. Documents longer than the context used in training are split into parts, as
described in Chapter~\ref{sec:setup}. Finally, the experiments target one class of hardware, a consumer laptop with
unified memory, and every system, including every baseline, runs on the same machine.

\section{Contributions}
\label{sec:intro-contributions}

The thesis makes four contributions.
\begin{enumerate}
\item \textbf{GroundOne}, an open specialist System One for sentence-level groundedness verification that runs entirely
on a laptop. Thanks to a marker placed after every sentence, it reads all the verdicts of an answer in one forward pass.
It is released as an installable package with a command-line interface and a local HTTP API, which also includes an
interface compatible with existing System One clients.
\item \textbf{A dataset built around the task} from public human annotations, with three test sets that were validated
without new annotation and frozen before any model was trained, and with leakage controlled at the level of documents.
\item \textbf{A controlled decomposition of the gains.} The comparison isolates the contribution of specialisation, of the
restricted readout and of model size. GroundOne is compared with cues that never verify a sentence, untrained models,
generative judges with and without probabilities, a general-purpose System One, a general-purpose decision model of the
same size and specialised encoders, all on the same hardware and with intervals obtained by cluster bootstrap.
\item \textbf{A study of calibration, time and energy}: how the probabilities can be calibrated and turned into consistent
labels, how much they help in deferring uncertain decisions, and how the time and energy of a verification grow with the
length of the answer. The study is set in an analysis of local inference that combines vendor specifications with own
measurements to bound, for each class of consumer hardware, what a verifier that generates text must pay and a System One
does not.
\end{enumerate}

\section{Main results}
\label{sec:intro-results}

The results show that specialisation works in the domain it was aimed at. On the RAGTruth test set GroundOne reaches a
three-way balanced accuracy of 0.769 with 0.5B parameters and 0.832 with 3B, against 0.37 and 0.44 for the same models
before training. It ranks supported sentences better than the specialised encoders, and it clearly outperforms both a
general-purpose System One eight times larger than the 0.5B and a general-purpose decision model of the same size as the
3B. On LLM-AggreFact, a benchmark drawn from different sources, the 3B is on par with the best small encoder.

The System One mechanism pays off in speed and in accuracy. When an answer contains fourteen or more sentences, GroundOne
verifies it about two times faster than a terse judge and six times faster than a verbose judge of the same size, with an
output that is always valid and with less energy. Restricting the decision to the three admissible letters is, by itself,
more accurate than letting the same model generate its verdict. After calibration, the probabilities of GroundOne are
close to the observed frequencies, whereas the probabilities of generative judges asked for the same output barely
separate correct from incorrect decisions.

Some aspects remain to be improved, and they point to natural directions for future work. Performance drops on domains
far from the training data, where the general-purpose decision model is stronger; the subtlest hallucinations, hidden in
sentences that otherwise follow the source faithfully, are detected only in part; and training is not yet fully stable
across random seeds.

\section{Thesis structure}
\label{sec:intro-structure}

Chapter~\ref{sec:local} is devoted to local inference: the reasons for keeping models on the device, open-weight and
small models, the cost model of inference on one device, the hardware landscape, runtimes and quantisation, energy, and
the consequence for verification. Chapter~\ref{sec:background} reviews System One models and the decision-model
ecosystem, hallucination in retrieval-augmented generation, verification methods and calibration. Chapter~\ref{sec:method} describes
GroundOne: the decision, the single-pass readout, the training objective, the calibration and the served output.
Chapter~\ref{ch:data} defines the labels and describes how the training and test sets were built, checked and frozen.
Chapter~\ref{sec:setup} presents the trained models, the baselines, the metrics and the measurement protocol.
Chapter~\ref{sec:results} reports the results, Chapter~\ref{sec:discussion} discusses what they mean and where they stop,
and Chapter~\ref{sec:conclusion} concludes. Appendix~\ref{app:examples} shows the output of every system on the example
answer described in Section~\ref{sec:intro-hallucination}.
